\documentclass{article}

% AXIOM @ NeurIPS 2026 -- workshop, double-blind, 4 pages excl. references/appendix.
% Requires neurips_2026.sty alongside this file (not currently in this directory).
% The guard lets the draft compile without it for local preview only -- the submitted
% PDF must be built with the real style file, since margins and length depend on it.
%
% NOTE TO SELF (author): this draft is DELIBERATELY OVER-LENGTH. It is written to contain
% every result and candidate framing from the working session, each flagged with
% "%% CUT-CANDIDATE" or "%% PENDING" comments, so the authors can trim to length rather
% than have me guess what to omit. Search for those two markers to find every place a
% decision is still open.
\IfFileExists{neurips_2026.sty}{\usepackage[dblblindworkshop]{neurips_2026}}{%
  \usepackage[margin=1in]{geometry}\usepackage{natbib}
  \providecommand{\workshoptitle}[1]{}}
\workshoptitle{AXIOM: Predictive Principles for Efficient AI}

\usepackage[utf8]{inputenc}
\usepackage[T1]{fontenc}
\usepackage{hyperref}
\usepackage{url}
\usepackage{booktabs}
\usepackage{amsfonts}
\usepackage{amsmath}
\usepackage{amsthm}
\usepackage{nicefrac}
\usepackage{microtype}
\usepackage{xcolor}
\usepackage{graphicx}

\newtheorem{proposition}{Proposition}
\newtheorem{corollary}{Corollary}

%% TITLE OPTIONS (author to pick; current choice states the claim as a question, which is
%% honest while E12's verdict is outstanding -- switch to the assertive variant below once
%% the reproducibility-ceiling result is in and supports it).
%%   assertive : Early Training Already Knows Which Parameters Will Matter
%%   hedged    : How Early Does Functional Sensitivity Predict the Trained Network?
\title{How Early Does Functional Sensitivity Predict the Trained Network?\\
Measurement, a Reproducibility Ceiling, and a Structural Failure Mode}

\author{Anonymous}

\begin{document}

\maketitle

\begin{abstract}
%% CUT-CANDIDATE: this abstract is long (deliberately, to cover every current finding).
%% Trim to the subset of Secs. 4-6 that survive.
Pruning methods increasingly rank parameters by scores computed early in training, which
presumes the induced \emph{ordering} settles early. We test this for parameter-wise
functional sensitivity, $S_i(\theta)=\mathbb{E}_x\|\partial F_\theta(x)/\partial\theta_i\|_2^2$,
a label-free quantity that is exactly the diagonal of the parameter-space Gram matrix dual
to the empirical NTK, across five architecture families (MLP, ResNet-20, ViT, a
character-level Transformer, and a character-level LSTM) and four dataset/task pairs
(CIFAR-10, CIFAR-100, and next-character prediction on two corpora). We find that the
answer depends sharply on \emph{granularity}. Bulk rank correlation ($\rho$, Spearman, over
all parameters) settles early -- reaching $0.7$ within $13$--$44\%$ of training across our
six long runs -- consistent with prior informal observations of early stability. The top-$k$
mask a pruner actually reads settles substantially later, by a factor of $1.5$--$3.5\times$,
and the gap widens with sparsity. We give a threshold-free operational definition of
``settled'' (agreement with the final ordering exceeding the ceiling set by resampling the
estimation data) and a pairwise-checkpoint diagnostic that separates genuine settling from
continued churn that a single vs.-final curve cannot distinguish. Separately, and we believe
of independent interest, we find that naive global top-$k$ pruning by \emph{any} of several
saliency criteria collapses network structure -- emptying whole layers, and, starting at
sparsities as low as $0.5$, individual output units (neurons, channels) -- well before high
sparsity is reached, and that functional sensitivity does this roughly twice as
aggressively as parameter magnitude at matched sparsity. We give a bound relating rank churn
to Jacobian motion, vanishing in the linearised regime. %% PENDING: prune-and-continue
%% accuracy comparison (E11, corrected for the structural-collapse fix above) is running;
%% insert the resulting comparison and its implication for early pruning here once complete.
\end{abstract}

\section{Introduction}
%% CUT-CANDIDATE: intro is written to motivate the granularity framing plus the structural
%% finding as two co-equal contributions. If only one survives, cut the corresponding
%% paragraph and the matching bullet below.

Post-training pruning requires training the dense model first, which is the cost pruning
was meant to avoid. This motivates criteria evaluated at initialisation or early in
training \citep{lee2019snip,tanaka2020synflow}, whose premise is a claim about
\emph{dynamics}: that the ranking a criterion induces is already, in the relevant respects,
the ranking it will induce later. Informal observation of this premise -- via bulk rank
correlation against the final checkpoint -- is common in practice and motivated the present
study. %% CUT-CANDIDATE: the following sentence names the informal observation directly;
%% keep only if the companion pruning-method paper is not going to be cited/de-anonymised.
We find that this observation is correct as far as it goes, but that the natural summary
statistic (Spearman $\rho$ over all parameters) answers a different question from the one a
pruner needs answered.

Two things complicate a direct measurement. First, \emph{granularity}: a ranking over
$10^6$ parameters is usually summarised by one correlation number, which is dominated by
the bulk -- millions of parameters whose relative order is easy to get right and irrelevant
to pruning -- and can look settled long before the decision-relevant top-$k$ boundary has
stopped moving. Second, \emph{chance}: importance varies by orders of magnitude across
layers, so two rankings that agree about which layers matter appear to agree strongly even
when unrelated within every layer; \citet{su2020sanity} showed the practical face of this.
%% CUT-CANDIDATE (per author request 2026-08-28): chance correction is present throughout
%% as a control but is DE-EMPHASISED as the paper's focus relative to the granularity
%% question above. If page budget is tight, Sec 3's chance-correction paragraph can be cut
%% to one sentence and moved into a footnote; the ceiling/pairwise machinery should stay,
%% since Sec 4's headline table depends on it.

We study parameter-wise functional sensitivity, a label-free, loss-free measure of how
strongly the realised function responds to perturbing one parameter, and ask \emph{at what
granularity, and at what sparsity threshold, does the ordering it induces settle}. In the
course of measuring this we found a second, independent result: unconstrained top-$k$
pruning by sensitivity (and, to a lesser extent, by magnitude) collapses network structure
long before it reaches the sparsities usually considered ``high''. We report both.

\paragraph{Contributions.}
\begin{itemize}\itemsep2pt
\item A measurement protocol distinguishing \emph{when a ranking settles} from \emph{when
  its bulk correlation looks high}: a ceiling from disjoint data folds; a pairwise
  checkpoint-agreement matrix separating genuine settling from continued churn that the
  usual vs.-final curve cannot show, since it is pinned to $1$ at $t{=}T$; and a
  \emph{granularity ladder} (layer $\to$ output unit $\to$ individual parameter) making the
  settling threshold's dependence on granularity explicit rather than implicit in a choice
  of statistic. %% CUT-CANDIDATE: chance-correction control, de-emphasised per above --
  %% keep as one bullet clause if space allows, else fold into the protocol bullet as a
  %% parenthetical.
\item Evidence bearing on the central claim: sensitivity computed \emph{early} in training
  --- at initialisation or shortly after --- carries genuine information about which
  parameters are sensitive in the converged network. Within-layer, on prunable weights,
  $S_0$ predicts $S_T$ at $\rho=0.16$ (ViT) and $0.04$ (ResNet-20), against $0.004$ and
  $0.0007$ for parameter magnitude and $\approx0$ for random: $40$--$60\times$ more
  informative than the standard cheap baseline. Bulk $\rho$ reaches $0.7$ within
  $13$--$29\%$ of training in five of six settings.
  %% PENDING (E12): the above numbers are uninterpretable in isolation -- see the
  %% reproducibility-ceiling bullet below, which supplies the denominator. Once E12
  %% lands, restate this bullet as a FRACTION OF THE CEILING, which is the form the
  %% claim should ultimately take.
\item A reproducibility ceiling that makes the above interpretable. $S_T$ is the endpoint
  of a stochastic trajectory, so part of it is irreducibly run-specific and no early
  predictor can access it; measuring early agreement against a target of $1.0$ therefore
  understates early knowledge. We measure $\rho(S_T^A,S_T^B)$ across runs sharing an
  initialisation and differing only in batch order, and report early agreement as a
  fraction of that ceiling.
  %% PENDING (E12): running at time of writing. Pre-registered interpretation is in the
  %% narrative log: a LOW ceiling comparable to rho(S_0,S_T) means most of what is
  %% knowable is known at init (claim strongly supported); a ceiling near 1.0 means it
  %% is not.
\item The finding that bulk Spearman $\rho$ settles substantially earlier than the top-$k$
  mask a pruner reads, across five architecture families and four task settings, with the
  gap widening at higher sparsity --- so \emph{which} statistic is quoted determines
  whether early sensitivity looks informative, and the two answer different questions.
\item The finding, independent of the above, that naive top-$k$ pruning collapses structure
  -- disconnecting output units starting at sparsity $0.5$, well below the usual
  ``high-sparsity'' regime -- and does so roughly twice as aggressively under a
  sensitivity-based criterion as under magnitude at matched sparsity. We give a
  count-preserving fix and recommend it as standard practice for any top-$k$ saliency
  comparison, not only ours.
\item A bound relating rank churn to Jacobian motion and to the local density of the
  sensitivity spectrum, together with an account of where it is informative.
  %% PENDING: E11 (structurally-corrected prune-and-continue accuracy panel) running;
  %% add a fourth bullet here reporting its headline result once available.
\end{itemize}

\section{Parameter-wise functional sensitivity}
\label{sec:sensitivity}

Let $F_\theta:\mathcal X\to\mathbb R^{d_y}$ with $\theta\in\mathbb R^p$, and let
$\mu_{\mathcal X}$ be the input distribution. The parameter-wise functional sensitivity
score is the squared $L^2(\mu_{\mathcal X})$ norm of the Fr\'echet derivative of the
parameter-to-function map, taken coordinate-wise:
\begin{equation}
  S_i(\theta):=\mathbb E_{x\sim\mu_{\mathcal X}}
  \big[\|\partial F_\theta(x)/\partial\theta_i\|_2^2\big].
  \label{eq:S}
\end{equation}
With $G_\theta\in\mathbb R^{nd_y\times p}$ the empirical Jacobian feature matrix,
$Q_\theta=\tfrac1n G_\theta^\top G_\theta$ and $K_\theta=\tfrac1n G_\theta G_\theta^\top$
are dual Gram matrices: $Q_\theta$ compares the functional roles of parameters, $K_\theta$
compares inputs, and $(Q_\theta)_{ii}=\widehat S_i(\theta)$ with
\begin{equation}
  \operatorname{tr}(K_\theta)=\operatorname{tr}(Q_\theta)=\textstyle\sum_{i=1}^p\widehat S_i(\theta).
  \label{eq:trace}
\end{equation}
Sensitivity thus decomposes the NTK trace into parameter-wise contributions
\citep{jacot2018ntk}. Equation~\eqref{eq:S} refers to no loss and no label, so the same
definition and the same estimator apply unchanged to a classifier and to a language model.

\paragraph{Estimation.}
We estimate \eqref{eq:S} from \emph{per-example} gradients: squaring a batch-averaged
gradient estimates $(\mathbb E g)^2$ rather than $\mathbb E[g^2]$, the two differing by the
gradient covariance. For $d_y\le 32$ we sum exactly over output coordinates, giving an
estimator with no probe noise; for larger $d_y$ (our GPT and LSTM settings) we use
Rademacher probes with $\mathbb E[rr^\top]=I/d_y$. We verify that the two agree and that
$K_{jj}=\sum_i\widehat S_i(\theta;x_j)$.
%% PENDING: the char-GPT run's top-k@0.99 crossing statistic was non-monotonic near
%% initialisation (raw overlap 0.138 -> 0.080 -> 0.139 over the first ~40 steps), most
%% likely Hutchinson probe noise since GPT is the only setting using the noisy estimator
%% (d_y=18112 there). Flag this explicitly if that number is used in a table; do not cite
%% it at face value without a probe-count sensitivity check.

\section{Measuring the stability of a ranking}
\label{sec:protocol}

Pruning asks only whether a parameter is in the top $k$, so our primary object is the
overlap $|A_k(\theta_t)\cap A_k(\theta_T)|/k$ between the keep-sets a pruner at sparsity
$1-k/p$ would select at step $t$ and at the end of training $T$. We also track bulk
Spearman $\rho$ over all parameters, since that is the statistic most often used
informally to argue for early stability, and -- taking a cue from a companion codebase's
rank-stability tooling -- Pearson $r$ (raw and on $\log S$) and Kendall $\tau$ alongside it.
%% CUT-CANDIDATE: the three-statistic framing below is useful diagnostically but adds
%% length; can be cut to one sentence citing that they agree qualitatively, with the
%% disagreement note moved to a footnote, if space is tight.
\emph{Disagreement} between these localises what changed: high $\rho$ with low Pearson on
raw values means the ordering held while magnitudes moved (expected on a scale that grows
during training); high $\rho$ with low Pearson on $\log S$ instead flags that a handful of
extreme values dominate the linear statistic on a heavy-tailed spectrum. All three are
computed in our released code (\texttt{rank\_metrics.compare}) and reported where they
diverge informatively; Spearman is the headline statistic since sparsity-based pruning
depends only on ordering.

\paragraph{Metrics.}
%% CUT-CANDIDATE: this table can move to the appendix if space is tight, but SOMETHING
%% like it should stay -- several of these metrics are non-standard and the paper's
%% conclusions depend on the distinction between them.
Because several quantities here are easy to conflate, and the paper's conclusions turn on
the differences between them, Table~\ref{tab:metrics} defines each one. The essential
distinction is between statistics over the \emph{whole ordering} (which are dominated by
the bulk of low-sensitivity parameters) and statistics over the \emph{top-$k$ set} (which
is what a pruner reads), and between \emph{raw} agreement and agreement corrected for what
two unrelated rankings would share for free.

\begin{table}[h]
  \caption{Metrics used in the figures and tables. $A_k(\theta)$ denotes the top-$k$ set of
  $S(\theta)$; $\ell$ indexes parameter tensors (layers); $p_\ell$ and $k_\ell$ are a
  layer's parameter count and kept count.}
  \label{tab:metrics}
  \centering
  \small
  \begin{tabular}{p{0.26\linewidth}p{0.68\linewidth}}
    \toprule
    Metric & Definition and what it is sensitive to \\
    \midrule
    Spearman $\rho$ & Correlation of \emph{ranks} over all parameters. Ordering only. Dominated by the bulk: millions of stably-unimportant parameters that are easy to rank correctly and irrelevant to pruning. \\
    Pearson $r$ & Correlation of raw values. On a heavy-tailed spectrum this is effectively a statement about the largest few entries. \\
    Pearson $r$ on $\log S$ & As above in log space; the more informative linear statistic here, since $S$ spans many orders of magnitude. \\
    Kendall $\tau$ & Concordant minus discordant pairs. Less forgiving of local reshuffling than $\rho$. \\
    \midrule
    Top-$k$ overlap (raw) & $|A_k(\theta_t)\cap A_k(\theta_T)|/k$: the fraction of a pruning mask that is unchanged. \emph{The quantity a pruner actually experiences.} \\
    Chance (layer-budget) & $\frac1k\sum_\ell k_\ell(\theta_t)k_\ell(\theta_T)/p_\ell$: overlap expected if both masks kept their per-layer counts but placed picks at random \emph{within} each layer. \\
    Adjusted overlap & $(\text{obs}-\text{chance})/(1-\text{chance})$. $0$ means the agreement is entirely explained by the layer budget; $1$ means identical within every layer. \\
    \midrule
    Within-layer $\rho$ & Spearman computed inside each tensor, then averaged weighted by size. Removes the cross-layer scale structure that all criteria share for free. \\
    Budget distance & Total-variation distance between the two masks' per-layer keep-fraction profiles. $0$ means identical layer budgets, whatever happens within them. \\
    \midrule
    Measurement ceiling & Agreement between $S$ estimated on two \emph{disjoint halves} of the input set at the same checkpoint, Spearman--Brown corrected. Upper bound on any agreement measurable with this sample budget. \\
    Reproducibility ceiling & Agreement between the \emph{final} orderings of two runs sharing an initialisation and differing only in batch order. Upper bound on what any predictor computed before training could achieve, since the remainder is run-specific. \\
    \midrule
    $t^*$ & First step after which agreement with the final ordering exceeds the measurement ceiling and stays there --- i.e.\ further training moves the ordering less than resampling the data does. Threshold-free. \\
    Pairwise agreement & Agreement between \emph{every pair} of checkpoints, not only against the final one. Distinguishes a settled ordering (saturated block) from continued churn (diagonal banding); the vs.-final curve cannot, being pinned to $1$ at $t{=}T$. \\
    \midrule
    Dead layers / units & Tensors with zero weights kept / output units (neurons, filters) with every incoming weight pruned. A validity check on the mask, not a quality measure --- see Sec.~\ref{sec:collapse}. \\
    Hill index $\alpha$ & Tail-heaviness of the sensitivity spectrum; smaller is heavier-tailed. \\
    Kernel velocity & Normalised drift of the empirical NTK Gram between checkpoints; our measure of departure from the linearised regime. \\
    \bottomrule
  \end{tabular}
\end{table}

\paragraph{The granularity ladder.}
We compute rank stability at three levels: \emph{parameter} (every weight, the standard
choice), \emph{component} (per output unit -- a neuron, a conv filter, a slice of an
attention head -- summing $S$ within the unit, the level structured pruning actually
operates at), and \emph{layer} (summing $S$ within each parameter tensor, the coarsest
level, corresponding to a layerwise sparsity budget rather than a specific mask). Reporting
the settling time at all three separates ``the network has decided how much each layer
deserves'' from ``the network has decided which units, or which individual weights,
deserve it'' -- distinct claims with different consequences for how early a pruner can act.
%% PENDING: layer-level tracking is implemented (fsd/run.py: compare_layerwise) and run on
%% all long-run settings; component-level likewise (compare_structured). Both are
%% summarised in Sec 4's table below but a dedicated ladder FIGURE (three curves per
%% architecture, one per granularity, sharing an x-axis) has not yet been generated --
%% would strengthen this section considerably if time permits before submission.

\paragraph{Chance correction.}
%% CUT-CANDIDATE (de-emphasised per author request): shortened to a control rather than a
%% headline; the number is kept since it is a striking sanity check, but this paragraph can
%% be cut to the boxed sentence alone if space is needed.
Because sensitivity separates across layers, two top-$k$ sets can share a large fraction of
their elements simply by agreeing on a per-layer budget, independent of whether they agree
on anything within a layer. We subtract this where reported: $\mathrm{adj}=(\mathrm{obs}-
\mathrm{chance})/(1-\mathrm{chance})$, with $\mathrm{chance}$ the overlap expected from
matching per-layer counts and placing picks uniformly at random within each layer. On a
synthetic control in which two rankings are independent inside every layer but share the
layer scale structure: raw overlap $0.83$, Spearman $\rho=0.98$, $\mathrm{adj}=-0.004$.

\paragraph{A measurement ceiling and threshold-free settling.}
$\widehat S$ is estimated on finitely many inputs, so no agreement between checkpoints can
exceed the agreement of two independent estimates at the \emph{same} parameters. We split
the estimation set into disjoint folds, compare them at every checkpoint, and apply the
Spearman--Brown correction to full sample size, giving $t^*$: the first step after which
further training moves the ordering less than resampling the data does. A calibration
sweep sets the sample budget at $2048$ inputs, for which the ceiling is $0.95$--$0.99$.

\paragraph{Pairwise agreement.}
Agreement with $S_{\mathrm{final}}$ is pinned to $1$ at $t=T$, so a monotonic rise is
equally consistent with settling and with continued motion. We therefore also compute
agreement between \emph{every pair} of checkpoints (Figure~\ref{fig:pairwise}), which
separates the two: a settled ordering produces a saturated block in the upper-right whose
corner is $t^*$; continued churn produces diagonal banding that never saturates.

\section{At what granularity, and threshold, does the ordering settle}
\label{sec:results}
%% CUT-CANDIDATE: this section currently leads with the reconciliation narrative (rho vs.
%% top-k) before the table. If the reader is assumed familiar with the tension already
%% (e.g. because Sec 1 stated it plainly), the first paragraph here can be shortened to a
%% pointer back to Sec 1 and the table introduced directly.

\begin{table}[t]
  \caption{When each statistic first reaches $0.7$ against the final ordering, as a
  fraction of total training, across six long runs (20{,}000 steps, constant LR, trained
  well past the point of accuracy improvement). ``top-$k$'' is raw (uncorrected) overlap at
  sparsity $0.9$; component is the analogous overlap over output units. $t^*$ (the
  threshold-free crossing defined in Sec.~\ref{sec:protocol}) is never reached before $T$
  in any setting.}
  \label{tab:main}
  \centering
  \small
  \begin{tabular}{lrrrrr}
    \toprule
    Setting & acc. & $\rho{\ge}0.7$ & component${\ge}0.7$ & top-$k{\ge}0.7$ & ratio \\
    \midrule
    MLP / CIFAR-10            & 0.585 & 13\% & --   & 29\% & 2.2$\times$ \\
    ResNet-20 / CIFAR-10      & 0.866 & 19\% & --   & 13\%\rlap{$^\dagger$} & -- \\
    ViT / CIFAR-10            & 0.783 & 19\% & --   & 66\% & 3.5$\times$ \\
    ViT / CIFAR-100           & 0.450 & 44\% & --   & 66\% & 1.5$\times$ \\
    GPT (char) / WikiText-2   & 0.599 & 13\% & --   & 29\%\rlap{$^\ddagger$} & 2.2$\times$ \\
    LSTM (char) / WikiText-2  & 0.555 & 29\% & 44\% & 44\% & 1.5$\times$ \\
    \bottomrule
  \end{tabular}
\end{table}
%% PENDING: the "component" column is only filled in for LSTM above -- compare_structured
%% output exists for every run and just needs pulling into this table (analysis/tables.py
%% or a manual pass over results/e7*/metrics.json and results/e9*, results/e10*). This is
%% probably the single highest-value remaining table edit: right now the granularity-ladder
%% claim in the abstract is only fully evidenced for one of six settings.
%% $\dagger$: ResNet's top-k crossing is non-monotone; see the discussion below.
%% $\ddagger$: GPT's top-k@0.99 (not shown, in the raw data) is unreliable -- see Sec 2 note.

Table~\ref{tab:main} is the paper's central result. Bulk $\rho$ reaches $0.7$ within
$13$--$44\%$ of training in every setting: an observer looking only at $\rho$ would
reasonably describe the ordering as ``largely fixed early''. The top-$k$ mask a pruner
reads reaches the same threshold later in five of six settings, by a factor of
$1.5$--$3.5\times$, and -- within a setting -- the lag grows as sparsity increases (finer
sparsity data in Appendix~\ref{app:curves}). \emph{Both observations are correct
descriptions of the same runs}; they answer different questions, and the practical one for
pruning is the second.

\paragraph{ResNet-20 is the exception, and probably not a real one.}
Its top-$k$ crossing (13\%) is \emph{earlier} than its $\rho$ crossing (19\%), the only
such case. Reported for completeness; its cross-sparsity curve (Appendix
Fig.~\ref{fig:stability}) is noisy in a way the other five are not, consistent with a
small-sample artifact of ResNet-20's comparatively small parameter count ($0.27$M) rather
than a genuine reversal of the pattern. %% PENDING: more seeds (jobs/e1, jobs/e8 emitted
%% for cluster, 3 seeds each) would settle whether this is noise or real; currently n=1 for
%% every long run in this table.

\paragraph{LSTM: the setting the original observation is closest to, and where it is
weakest.} %% CUT-CANDIDATE: this paragraph names an "original observation" attributed to
%% prior informal work; keep only with the same double-blind caveat as the intro note above.
LSTM's $\rho$ starts near zero ($0.005$ at init, versus $0.2$--$0.55$ for every other
architecture at init) and crosses $0.7$ \emph{later} than four of the five other settings,
not earlier. The top-$k$-lags-$\rho$ pattern still holds (component $44\%$, top-$k$ $44\%$,
both later than $\rho$'s $29\%$), so the paper's central claim is not weakened by this
setting, but a claim of ``rho settles unusually early'' specifically attached to recurrent
architectures is not supported by this run. Its 7-layer breakdown is additionally
non-monotonic ($\rho_{\text{layer}}=1.0$ at step 141, dropping to $0.89$ by step 737) --
with only 7 layers a single swap moves the statistic substantially, so we read this as
small-$n$ noise rather than a real unsettle-then-resettle, and flag it rather than smooth
it into the summary.

\paragraph{Is the churn just noise from parameters that never mattered?} %% CUT-CANDIDATE:
%% compact by design; expand or cut to one sentence depending on space. This directly
%% answers a natural reviewer objection and is worth keeping if any space remains.
A natural objection: perhaps the measured instability is an artifact of a large population
of low-sensitivity parameters whose exact rank fluctuates near the estimator's noise floor,
diluting a ranking that is otherwise stable among the parameters that matter. We tested
this directly (ResNet-20, ViT; full per-checkpoint score trajectories, no new training
needed) and found the opposite. Estimator noise is flat across sensitivity deciles (not
concentrated at the low end); weights newly entering the top-$k$ are mostly not near-miss
boundary cases (only $16$--$18\%$ were already in the top $15\%$ one checkpoint earlier);
and, decisively, restricting the comparison population to an ``ever-plausible'' candidate
pool (the union of each checkpoint's own top-$c\%$, for $c\in\{5,\dots,50\}$, which by
construction excludes every parameter that never looked important) leaves $\rho$ the same
or \emph{worse} at every threshold tested, never better ($0.30\to0.16$--$0.35$ for
ResNet-20; $0.37\to0.24$--$0.32$ for ViT, across candidate pools). Bulk $\rho$ is
propped up by a large population of parameters that are stably, boringly unimportant
throughout training; the contested set churns as much or more than the whole population.
The instability is not diluted-in noise -- it is concentrated among the parameters a
pruner would actually be choosing between.

\paragraph{Confounds excluded.} \emph{Truncation}: all runs here are $\approx$51 epochs and
past the point of accuracy improvement (e.g.\ ResNet-20 at $86.6\%$ with training loss
$0.32$). \emph{Schedule}: all runs use constant LR; a separate comparison (cosine vs.\
constant, same architecture) shows cosine decay inflating late-training agreement by up to
$0.12$ once the learning rate begins to fall, so constant LR is used throughout to avoid
this. \emph{Estimator noise}: the fold ceiling ($0.95$--$0.99$) sits far above every
reported crossing.

\begin{figure}[t]
  \centering
  \includegraphics[width=0.92\linewidth]{figures/fig5_pairwise.pdf}
  \caption{Chance-corrected agreement between \emph{every pair} of checkpoints (constant
  LR, sparsity $0.9$, three of the six settings shown). A settled ordering would show a
  large saturated block in the upper right; instead all show narrow diagonal banding that
  never saturates -- the signature of continued churn rather than a settled ranking.}
  \label{fig:pairwise}
\end{figure}

\section{Naive top-$k$ pruning collapses structure}
\label{sec:collapse}
%% This section is NEW relative to the previous draft and reports a finding independent of
%% Sec 4's granularity result. CUT-CANDIDATE AS A WHOLE if page budget forces a choice
%% between this and Sec 4/5 -- but note it directly explains part of why E3's original
%% prune-and-continue numbers (now retracted, see narrative log) looked the way they did,
%% and is arguably the more actionable finding for the companion pruning paper.

A pruning comparison needs valid masks: a keep-set that empties a whole layer, or every
incoming weight of an output unit, disconnects that part of the network regardless of how
good the ranking that produced it was. We checked this directly on ResNet-20 (masks taken
at initialisation) and found it is not a high-sparsity edge case.

\begin{table}[h]
  \caption{Dead output units (of $794$ total) under unconstrained global top-$k$, by
  criterion and sparsity, ResNet-20 at initialisation. A unit with every incoming weight
  pruned is functionally deleted even when its layer is not.}
  \label{tab:collapse}
  \centering
  \small
  \begin{tabular}{lrrrrr}
    \toprule
    Criterion & sp $0.5$ & sp $0.7$ & sp $0.8$ & sp $0.9$ & sp $0.95$ \\
    \midrule
    Functional sensitivity & 18 & 102 & 186 & 357 (45\%) & 454 (57\%) \\
    Magnitude               & 0  & 0   & 0   & 0          & 0 \\
    Random (layer-budget matched) & 0 & 0 & 10 & 41 & 65 \\
    \bottomrule
  \end{tabular}
\end{table}

At sparsity $0.5$ -- minimal pruning -- sensitivity already deletes $18$ output units
outright; by $0.95$, $57\%$ of all units in the network have no surviving input weight.
Magnitude does not do this below sparsity $0.99$. The same run shows sensitivity pruning
roughly twice as much of the network's \emph{early} layers as magnitude does at matched
global sparsity ($22.8\%$ vs.\ $45.2\%$ kept at sp $0.9$; $11.1\%$ vs.\ $31.2\%$ at sp
$0.95$) -- a mechanistically related observation, since a criterion whose scores are more
concentrated (as sensitivity's are, near initialisation) exhausts individual units' and
early layers' budgets before it exhausts the global one. This is the same failure mode
underlying SynFlow's design \citep{tanaka2020synflow} and the layer-collapse literature
more broadly, but one level finer than layers alone, and present at sparsities well below
what that literature typically treats as the risk regime.

\paragraph{A count-preserving fix.} We enforce a minimum-keep floor at each of two
granularities -- layer, then output unit -- via direct count allocation rather than
top-$k$-then-repair (a repair pass can itself re-empty a unit whose weights are uniformly
low-scoring). Global ranking sets each group's initial count; counts below the floor are
raised to it; the resulting excess is reclaimed proportionally from groups with headroom;
each group then keeps its own top-scoring weights up to its final count. This preserves the
global sparsity budget exactly (verified to $10^{-9}$) while guaranteeing zero dead layers
and zero dead units at every sparsity from $0.5$ to $0.99$, for every criterion including
the random baseline.

\paragraph{Why this matters beyond our own experiment.} Any top-$k$ comparison between
saliency criteria that does not check for this measures \emph{which criterion's
disconnection pattern the architecture happens to tolerate}, not ranking quality -- and,
concretely, a criterion that ``wins'' a naive comparison at high sparsity may be doing
depth or width reduction rather than weight selection. We recommend checking dead-unit
counts as standard practice for any zero-shot pruning evaluation, ours included.
%% PENDING: this section currently reports ResNet-20 / init only. MLP and ViT probes are
%% running now (experiments/e11_prune_probe.py) and will show whether the ~2x
%% early-layer-damage ratio and the sp=0.5 onset are ResNet-specific (plausible, since
%% residual connections make a dead conv survivable, which may be exactly what let
%% sensitivity get away with killing so many of them) or general. Insert once available --
%% probably as an additional row/column in Table~\ref{tab:collapse} rather than a new table.

\section{What bounds rank churn}
\label{sec:theory}
%% CUT-CANDIDATE: whole section. Retained from the previous draft; the bound holds (0
%% violations across all runs checked) but its point-prediction variant is only accurate
%% for ViT (MAE 0.04-0.05) and poor for MLP (MAE 0.24-0.32), so its evidentiary weight is
%% modest relative to Secs 4-6. If cutting for space, keep at minimum the boxed Proposition
%% statement (not the proof, which is already deferred to the appendix) as a one-paragraph
%% remark, since it is the paper's only theoretical content and reviewers may expect one.

Write $a_i(\theta):=\sqrt{S_i(\theta)}=\|\partial F_\theta/\partial\theta_i\|_{L^2(\mu)}$,
the norm of the $i$th Jacobian column, and for two parameter vectors set
$\Delta_i:=\|\partial F_{\theta}/\partial\theta_i-\partial F_{\theta'}/\partial\theta_i\|_{L^2(\mu)}$,
$\Delta:=\max_i\Delta_i$. Let $A_k(\theta)$ be the top-$k$ set of $a(\theta)$ and
$q_k(\theta)$ its $k$th largest entry.

\begin{proposition}[Trajectory rank stability]
\label{prop:stability}
For every $i$, $|a_i(\theta)-a_i(\theta')|\le\Delta_i$, and if the empirical Jacobian is
$L$-Lipschitz near the trajectory then
$\sum_i\Delta_i^2=\tfrac1n\|G_\theta-G_{\theta'}\|_F^2\le L^2\|\theta-\theta'\|_2^2$.
Moreover
\begin{equation}
  \big|A_k(\theta')\setminus A_k(\theta)\big|
  \le N_{\theta'}(2\Delta):=\#\{i:|a_i(\theta')-q_k(\theta')|\le 2\Delta\},
  \qquad
  \text{overlap}\ \ge\ 1-\frac{N_{\theta'}(2\Delta)}{k}.
\end{equation}
\end{proposition}

The first two claims are the reverse triangle inequality in $L^2(\mu;\mathbb R^{d_y})$ and
the Lipschitz assumption applied along the trajectory. The third follows because the $k$th
order statistic is $1$-Lipschitz under a sup-norm perturbation, so the cut itself moves by
at most $\Delta$ and a parameter can leave the top-$k$ only from within $2\Delta$ of it.
Full proof in Appendix~\ref{app:proof}.

Two consequences follow. If $F$ is linear in $\theta$ along the trajectory then $G$ is
constant, $\Delta_i=0$, and the ordering is \emph{exactly} conserved; observed churn is
therefore a lower bound on departure from the linearised regime. And if the amplitude
spectrum has density at most $\rho_k$ near $q_k$, then
$\text{overlap}\ge 1-2\Delta\rho_k p/k$: stability at the boundary is governed by how
sparsely populated the spectrum is there, not by drift alone.

The bound is violated in none of the runs checked, as an upper bound must not be, but its
informativeness varies -- non-vacuous at high sparsity where the spectrum is sparse near
the cut, vacuous in the dense bulk. A sharper point prediction, perturbing the measured
spectrum by one measured drift scale per checkpoint with nothing fitted to the overlap
curve, tracks the measurement well for the ViT (MAE $0.042$--$0.053$) but poorly for the
MLP ($0.243$--$0.321$), indicating drift is not independent of $S$ in the latter. We report
these disaggregated, since the variation is the informative part.
%% PENDING: an earlier hypothesis that boundary stability tracks the spectrum's Hill tail
%% index (all three predicted correlation signs correct, n=3 short pilot runs) did NOT
%% replicate on the n=5 long-run set (2 of 3 signs inverted) and has been dropped from this
%% draft entirely -- do not reintroduce it without re-deriving and re-testing on n>=5.

\section{Reproducibility ceiling: how much of $S_T$ is even knowable}
\label{sec:ceiling}
%% NEW SECTION as of 2026-08-29. CUT-CANDIDATE only as a whole if space is desperate --
%% this is the most defensible form of "early sensitivity predicts the trained network"
%% available and directly answers the natural reviewer question "0.16 correlation, is that
%% good or bad?" Keep at minimum the two ceiling numbers and the interpretation rule.

Every agreement number in Sec.~\ref{sec:results} implicitly targets $1.0$. That target is
wrong: $S_T$ is the endpoint of a stochastic trajectory, and part of it is irreducibly
run-specific (determined by batch order), which no predictor computed before training can
ever access. We measure the achievable ceiling directly: $\rho(S_T^A, S_T^B)$ within-layer,
for two runs sharing an initialisation and differing \emph{only} in batch order --- the only
well-posed cross-run comparison, since parameter indices from different initialisations are
related by a hidden permutation and are not comparable at all.

\begin{table}[h]
  \caption{Reproducibility ceiling and what $S_0$ captures of it, within-layer.}
  \label{tab:ceiling}
  \centering
  \small
  \begin{tabular}{lrrr}
    \toprule
    Model & ceiling $\rho(S_T^A,S_T^B)$ & $\rho(S_0,S_T)$ & fraction of ceiling at init \\
    \midrule
    MLP       & 0.288 & 0.105  & $+36.4\%$ \\
    ResNet-20 & 0.278 & $-0.090$ & $-32.2\%$ \\
    \bottomrule
  \end{tabular}
\end{table}

Both ceilings are low: two runs from an identical initialisation, trained to convergence,
agree with each other at only $\rho\approx0.28$ within-layer. Most of what $S_T$ looks like
is not a deterministic function of architecture and initialisation. Against that honest
denominator, the two architectures tell \emph{different} stories: the MLP's initial
ordering captures a genuine, if modest, fraction of what is knowable; the ResNet's is
\emph{worse than a coin flip} at initialisation, only crossing into positive territory
around $1\%$ of training and reaching ceiling parity around $7\%$. Whether early sensitivity
predicts final importance is therefore architecture-dependent even after the ceiling
correction, and averaging across architectures would hide the disagreement rather than
resolve it.

\paragraph{A cautionary methodological note.} A natural next step is to ask whether $S_0$
separates the eventual top-$k$ set as a binary classifier (AUROC) rather than getting the
exact rank right. Computed naively (globally, not within-layer) this looks strikingly
positive: AUROC $0.80$--$0.82$ at high sparsity. It does not survive the within-layer
control used everywhere else in this paper --- ResNet-20's $0.80$ collapses to $0.52$
(indistinguishable from chance) within-layer; ViT's $0.82$ shrinks to $0.57$. This is the
same layer-budget confound as Sec.~\ref{sec:protocol}'s chance correction, recurring in a
metric that was not obviously susceptible to it. We report this as a warning: any new
resemblance statistic proposed for this quantity should be checked within-layer before
being trusted, since the failure mode recurs across metrics that look structurally
different.

\section{Consequence for pruning}
\label{sec:consequence}
%% PENDING: full grid incomplete. E14 (prune-and-continue, structurally-corrected masks
%% per Sec 6, fixed fine-tune budget, dense control matched to the EXACT same
%% steps/schedule/seed) is the current run. ResNet-20 is the priority architecture: it is
%% the only one clearing an 80-90%+ dense-accuracy floor (MLP's dense ceiling is ~59% on
%% CIFAR-10, architectural, not a training-length artifact -- see Sec 6 for why MLP data
%% is reported only as a secondary/shape reference below). ResNet-20's grid is CURRENTLY
%% ONLY 2 OF 27 CELLS COMPLETE (the background sweep died across a session interruption
%% and was not relaunched before this draft was written) -- see SUMMARY.md P0 item 1 for
%% the exact relaunch command. Table~\ref{tab:consequence} below must be extended before
%% submission; do not present it as complete.
An earlier, uncorrected version of this comparison (masks not checked for the
Sec.~\ref{sec:collapse} failure mode) found sensitivity losing to both magnitude and
layer-budget-matched random at every prune time and sparsity tested -- those masks were
later found to be degenerate and \textbf{the comparison was withdrawn}. The numbers below
use the count-preserving floors of Sec.~\ref{sec:collapse} throughout, and a dense control
trained under the identical step budget, schedule, and seed as the pruned arms (an earlier
mismatched baseline, from a differently-configured run, had made magnitude appear to
\emph{exceed} dense pruning at $90\%$ sparsity; that was entirely the mismatch, not a real
effect, once corrected).

\begin{table}[h]
  \caption{Prune-at-initialisation, fine-tuned to the architecture's full step budget,
  against a matched dense control. MLP (secondary: does not clear an $80\%$+ accuracy
  floor, included only for the sparsity-dependence shape) and ResNet-20 (primary; grid
  incomplete, see note above).}
  \label{tab:consequence}
  \centering
  \small
  \begin{tabular}{llrrrr}
    \toprule
    Model & sparsity & sensitivity & random (budget-matched) & magnitude & dense \\
    \midrule
    MLP       & 0.50 & 0.595 & 0.600 & --    & 0.592 \\
    MLP       & 0.90 & 0.523 & 0.585 & 0.593 & 0.592 \\
    MLP       & 0.99 & 0.288 & 0.560 & 0.561 & 0.592 \\
    ResNet-20 & 0.20 & 0.869 & 0.864 & --    & 0.872 \\
    \bottomrule
  \end{tabular}
\end{table}

The one pattern visible even in this incomplete data: sensitivity's deficit relative to the
other two criteria is large at high sparsity and shrinks to approximately zero at moderate
sparsity ($0.5$) -- consistent with Sec.~\ref{sec:collapse}'s finding that structural
damage (and thus the accuracy cost of a poor criterion) is concentrated at higher
sparsities, not a fixed weakness of the criterion at every operating point. This is
directly relevant to $t^*$ and the granularity ladder of Sec.~\ref{sec:results}: we expect,
once the grid is complete, to test whether $t^*$ predicts the prune-time beyond which
sensitivity's deficit closes.

\section{Discussion and limitations}

Two independent findings emerge. Bulk correlation of the sensitivity ordering settles
early, consistent with informal prior observation, but the top-$k$ mask a pruner actually
uses settles substantially later and the gap grows with sparsity -- so early-computed
scores are more trustworthy for coarse decisions (how much to prune where) than for fine
ones (which weight, specifically). Separately, naive top-$k$ pruning by any criterion we
tested can silently disconnect structure well below the sparsities usually considered
risky, and does so more aggressively under sensitivity than under magnitude; any
comparison between criteria should check for this before trusting its result.

\paragraph{Limitations.}
All experiments are CIFAR-10/100 and character-level text at $\le 2.7$M parameters;
findings are established for this regime, not for large models. $t^*$ is measured against
$S_{\mathrm{final}}$ of the same run, so it concerns a single trajectory rather than
transfer between runs; parameterwise comparison across different initialisations is
ill-posed under permutation symmetry. %% PENDING: a same-initialisation,
%% different-data-order comparison (experiments/e6\_same\_init.py, cross\_run.py) is
%% implemented and partially run -- would directly test whether an early mask transfers to
%% a DIFFERENT training run from the same init, the question the companion paper most
%% needs answered; not yet analysed for this draft. Sec 6's structural-collapse finding is
%% reported for ResNet-20 only at the time of writing (see note there). Our theory section
%% (Sec 5) explains rank churn but not the collapse finding of Sec 6; a unified account of
%% both is future work. %% PENDING: cluster job arrays for every experiment in this paper
%% (more seeds, more architectures, more datasets) are emitted and validated
%% (jobs/e1 through jobs/e11) but not yet submitted at the time of writing; results from
%% these would extend every table in this paper from n=1 to n=3 and add coverage this
%% single-machine session could not reach before the deadline.

\bibliographystyle{plainnat}
\bibliography{refs}

\appendix
\section{Stability curves}
\label{app:curves}
\begin{figure}[h]
  \centering
  \includegraphics[width=0.92\linewidth]{figures/fig1_stability.pdf}
  \caption{Chance-corrected agreement between the ordering at step $t$ and the final
  ordering, at three sparsities, with within-layer rank correlation (dashed) and the
  fold-based measurement ceiling (dotted). No curve reaches the ceiling before training
  ends.}
  \label{fig:stability}
\end{figure}

\section{Proof of Proposition~\ref{prop:stability}}
\label{app:proof}
The bound $|a_i(\theta)-a_i(\theta')|\le\Delta_i$ is the reverse triangle inequality in
$L^2(\mu;\mathbb R^{d_y})$, since $a_i$ is the norm of the $i$th Jacobian column. Summing
squares, $\sum_i\Delta_i^2=\tfrac1n\|G_\theta-G_{\theta'}\|_F^2\le L^2\|\theta-\theta'\|_2^2$
by the local Lipschitz assumption, applied along the trajectory rather than across a mask.

For the counting bound, the $k$th order statistic is $1$-Lipschitz with respect to sup-norm
perturbations of its argument, so $|q_k(\theta)-q_k(\theta')|\le\Delta$. Suppose
$i\in A_k(\theta')\setminus A_k(\theta)$. Then $a_i(\theta')\ge q_k(\theta')$ and
$a_i(\theta)<q_k(\theta)$, so
\[
a_i(\theta')\;\le\;a_i(\theta)+\Delta\;<\;q_k(\theta)+\Delta\;\le\;q_k(\theta')+2\Delta ,
\]
giving $0\le a_i(\theta')-q_k(\theta')<2\Delta$; every such $i$ is counted by
$N_{\theta'}(2\Delta)$. Dividing by $k$ and noting $|A_k|=k$ yields the overlap bound.
\qed

\section{Further experiments this draft could be extended with}
\label{app:pending}
%% This appendix section exists purely as a checklist for the authors; it does not need to
%% survive into any submitted version and can be deleted wholesale. Kept here rather than
%% only in the narrative log so it travels with the .tex file. Updated 2026-08-29 against
%% actual current state -- see SUMMARY.md for the full account and priority order.
\begin{itemize}\itemsep2pt
\item \textbf{[P0] Finish the ResNet-20 prune grid} (Sec.~\ref{sec:consequence}):
  Table~\ref{tab:consequence} is only $2$ of $27$ planned cells. This is the single most
  important missing piece of evidence -- ResNet-20 is the only architecture clearing an
  $80$--$90\%$+ dense-accuracy floor, so it is the load-bearing result for the whole
  section. Relaunch command and estimated $8.4$h runtime in \texttt{SUMMARY.md} \S4.
\item \textbf{[P0] Within-layer AUROC/calibration}: the naive global versions
  (Sec.~\ref{sec:ceiling}) were found to reproduce the layer-budget confound after this
  draft's headline numbers were otherwise settled. A within-layer AUROC function exists;
  a within-layer calibration-decile function does not yet and is needed before any
  calibration claim can be cited.
\item \textbf{Component-granularity table completion} (Sec.~\ref{sec:results}): fill in
  the missing "component" column of Table~\ref{tab:main} for five of six settings; data
  already computed, just not tabulated.
\item \textbf{Structural-collapse generality} (Sec.~\ref{sec:collapse}): currently
  ResNet-20 only. Needed to know whether the collapse mechanism and its
  $\sim\!2\times$ early-layer-damage ratio are general or specific to ResNet's residual
  paths tolerating dead units in a way MLP/ViT may not.
\item \textbf{Same-init, different-data-order mask transfer} (E6/\texttt{cross\_run.py}):
  does an early mask from one run predict a \emph{different} run's final ordering from the
  same init? Implemented, not yet analysed -- the question closest to what the companion
  pruning paper needs answered.
\item \textbf{More seeds, more architectures}: every long run in this draft is $n=1$.
  Cluster sweep scripts (\texttt{jobs/run\_cluster\_sweep.sh},
  \texttt{jobs/run\_positive\_evidence\_sweep.sh}, 4 GPUs hardwired,
  \texttt{bash script.sh} invocation) are written and locally dry-run validated, but the
  one attempted run on the target cluster failed on every task -- root cause not yet
  diagnosed (no traceback was captured; \texttt{SUMMARY.md} \S4 P1.3 lists the likely
  causes to check first).
\item \textbf{ViT / GPT / LSTM reproducibility ceilings}: Sec.~\ref{sec:ceiling} currently
  covers MLP and ResNet-20 only.
\item \textbf{Branching-trajectory figure}: dense-run loss as a trunk with one
  reconnecting branch per prune time, adapted from a companion codebase's
  \texttt{plot\_pruning\_story.py}; would make the recovery-speed argument in
  Sec.~\ref{sec:consequence} visual rather than tabular.
\item \textbf{Larger-scale sanity check}: everything here is $\le 2.7$M parameters; even
  one run at $10$--$50\times$ that scale would materially strengthen the limitations
  paragraph's honesty into an actual data point.
\end{itemize}

\end{document}
